% Nouveau EGA V (cote 25): archive cote 25, site page 52 (PDF page 53).
% typescript. Draft compared: none.
% AI-transcribed by Claude Opus 5.5 (claude-opus-5-5, adaptive thinking, effort high), Message Batches API, batch ega25b (msgbatch_01A7n5YuXubUPNK1zsepi8RG), 2026-11-13; not yet checked by eye.
\scanpage{25/52}
\mspage{47}

\xout{En vertu de 5.6, l'ensemble U des P tels que Y soit géom}

Supposons d'abord X géom normale. On sait déjà par 8.14. a)  que
l'en \xout{ex} peut trouver un fermé Z' de P de codi\xout{\unsure{mensi}}on $\geq$2, tel que
\hand{$\zeta\in$} P - Z'  implique que  $Y_{\hand{$\zeta$}}$  est  \uline{séparable} sur k(\hand{$\zeta$}) de  dim \inserted{d-}1; pour
un tel \hand{$\zeta$} , il revient donc au même que $Y_{\hand{$\zeta$}}$ soit géom irréd, ou géom
intègre, et  les deux problèmes  (respé et non respé) examinés dans
8.18. sont donc équivalents\unsure{.|!} D'autre part, en vertu de 5.6. l'ensemble
\inserted{U} des \hand{$\zeta\in$} $\hand{$\check{}$}$P  tels que  $Y_{\hand{$\zeta$}}$.  soit géom intègre de dim \hand{d-1} (l'hypothèse
dimensionnelle exprimant que $\phi_{\hand{$\zeta$}}$ est $X_{k(\hand{$\zeta$})}$-régulier ) est \uline{ouvert}.
Nous allons exhiber un ensemble ouvert \inserted[typed]{non vide,} évident  P - Z  cont\unsure{è|e}nu dans U,
en\xout{p} prenant pour  Z la réunion de  $\overline{g(\underline{Y}'^{\,\mathrm{sing}})}$  et des hyperplans $\check{H}_x$
de $\check{P}$  définis par les  f(x), \hand{$x\in$} T. \xout{Pour} Pour \hand{$\zeta\in$} $\check{P}$ - Z,  $Y_{\hand{$\zeta$}}$  est \uline{lisse}
de dimension \inserted{d-}1, et comme il est géom connexe par 6.1. , il est géom
intègre.  Il faut donc  exprimer qu\xout{\unsure{e chaque}} \inserted[typed]{toute} composante irréductible
de \inserted[typed]{co}dimension  1 de Z  \xout{\unsure{ne ren}}\inserted{\unsure{ren}}contre l'ouvert U.  Or ces composantes
irréductibles sont les  \hand{$\check{H}_x$}  (répetés éventuellement, peu importe), et
d'autre part $\overline{g(\underline{Y}'^{\,\mathrm{sing}})}$  lorsque  ce dernier  est bien de codim 1, i.e.
X' ``non réglée'' pour f, (NB on utilise l'irréductibilité de  $\underline{Y}'^{\,\mathrm{sing}}$ ).
D'ailleurs, pour que ce dernier ensemble rencontre  l'ouvert U,
il f et s  que  g($\underline{Y}'^{\,\mathrm{sing}}$)  (qui en contient un ouvert dense) rencontre
U, \xout{\unsure{i.e. qu'il existe x' X'(k) et une}} Cela prouve 8.18. dans ce
cas.  Lorsqu'on ne suppose pas X normal, on applique le résultat précé-%
dent au normalisé de X; le raisonnement est immédiat et je ne le détaille
pas ici.  \margin{\uline{NB} dans le cas \inserted{\unsure{bas}}respé, 8.18. est contenu dans 8.16. plus précis sauf pour d=2. C'est pour le cas non respé qu'il peut être utile de ne pas exiger d=2, mais seulement $d\geq 2$ ....}

Il reste à expliciter  les conditions a) et b) de 8.18. Cela con-%
duit à examiner de façon générale la situation suivante:  On suppose
X \deleted{intègre}\inserted{géom irréd,} sur k \deleted{alg clos}, on se donne une sous-variété linéaire  L \hand{de}

\margin{52}
